\section{Introduction}\label{sec:intro}

% ---- P1. The problem -------------------------------------------------------
Firms are adopting artificial intelligence (AI) rapidly, and early experimental and field evidence documents sizable productivity gains in specific tasks \citep{noy2023experimental,brynjolfsson2025generative}.
If these gains are labor-saving at the firm level, the first-order prediction is that incumbents will produce the same output with fewer workers \citep{acemoglu2019automation}.
The concern is sharpest for large firms, which have the data, capital, and organizational capacity to adopt early: in the European Union in 2025, 55\% of enterprises with 250 or more employees used AI, against 17\% of those with 10 to 49 employees \citep{eurostat2026isoc}.
This paper is about the second-order response.
A technology that makes existing firms more productive may also make new firms cheaper to create.
We ask whether, when AI reduces headcount at large incumbents, it also lowers the barriers to firm formation enough that employment in new small and mid-sized firms offsets the displacement.

% ---- P2. Why it matters ----------------------------------------------------
The answer matters for how economists and policymakers read the labor-market consequences of AI.
Most existing assessments measure the exposure of occupations and tasks \citep{eloundou2024gpts,acemoglu2025simple}, implicitly holding the set of firms fixed.
Yet net job creation in the United States is driven disproportionately by young firms rather than by small firms as such \citep{haltiwanger2013who}, and the long decline in business dynamism has been linked to slower reallocation and weaker productivity growth \citep{decker2014role,decker2020changing}.
If AI moves the entry margin, changes in the number and size of firms could rival the within-firm displacement that task-based accounts emphasize.
The question is equally pressing in Europe, where firm size distributions, entry costs, and labor-market institutions differ from those in the United States \citep{bartelsman2013cross}, and where a new regulatory framework for AI is being phased in \citep{euaiact2024}.

% ---- P3. Why it is hard ----------------------------------------------------
Answering the question is hard for three reasons.
First, AI is still diffusing, so the post-shock firm size distribution is not yet observed and any direct test is necessarily partial.
% [claim] speed asymmetry is plausible but should be supported (e.g. entrant scale-up horizons in BDS).
Second, the two channels operate at different speeds: incumbents can cut headcount within a budget cycle, whereas the entry, survival, and scale-up of new firms unfold over years, so short-run data will tend to overstate net displacement.
Third, the natural historical comparison is not clean: the cloud era coincided with a secular decline in the aggregate US startup rate \citep{decker2016declining,akcigit2021ten}, so any entry effect of cloud computing must be identified within exposed sectors and against a falling baseline rather than read off aggregate series.
A naive before-and-after comparison of aggregate employment therefore conflates the two channels and confounds both with unrelated trends.

% ---- P4. The gap and how this paper differs ---------------------------------
Prior work has studied each channel largely in isolation.
A large literature on automation studies displacement and reinstatement of labor within a given set of firms or industries \citep{autor2003skill,autor2015why,acemoglu2018race}, and work on superstar firms documents rising concentration alongside a falling labor share \citep{autor2020fall}.
Research on information technology shows that its effects depend on how firms are organized \citep{bloom2012americans,bloom2014distinct}, and recent evidence links AI investment to firm growth and product innovation \citep{babina2024artificial}.
Separately, research on cloud computing shows that cheaper experimentation changed how startups are financed \citep{ewens2018cost} and affected the survival and performance of young establishments \citep{jin2017economies}.
Canonical models of firm dynamics connect entry costs to the firm size distribution \citep{hopenhayn1992entry,luttmer2007selection}.
% [claim] "to our knowledge" requires a completed literature search before submission.
To our knowledge, these models have not been used to ask whether a shock that combines labor-saving technical change with lower entry costs raises or lowers aggregate employment.
We bring these strands together and use cloud computing as a disciplined analogue for the entry-cost component of the AI shock.

% ---- P5. What we do, what we find, and the limits --------------------------
We proceed in four steps.
We first set out a model of firm dynamics with endogenous entry and exit (\cref{sec:theory}).
We then use the cloud era of 2006--2020 to ask which parameter configurations reproduce what happened: lower computing costs coincided with rising employment concentration and a falling startup rate, a pattern the model matches only when a labor-complementary technology is combined with entry headwinds (\cref{sec:cloud}).
Next, we formalize the AI shock as a combination of displacement at incumbents and lower fixed costs of entry and operation (\cref{sec:ai-model}), and we calibrate the model and simulate it by Monte Carlo (\cref{sec:simulation}).
We find that the probability that AI raises net employment is 81.7\%, but the result is conditional: whether AI is labor-complementary ($\sigma_{\text{ai}} < 2.78$) or labor-displacing explains 71\% of the variance, while entry-cost reduction---the mechanism most analogous to cloud computing---explains only 3\%.
Below the displacement threshold, offset is nearly universal; above it, offset depends on whether AI-native entrants are substantially more productive than they would have been absent AI.
Finally, we derive empirical signatures that can be tested as AI diffuses (\cref{sec:signatures}).
Our analysis has clear limits: the analogy to cloud computing is imperfect, and because AI can substitute for labor at entrants as well as at incumbents, more firms need not mean more jobs.
In the discussion (\cref{sec:discussion}) we distinguish AI infrastructure providers (labs, hyperscalers) from AI application companies and document that the leading application firms are far more employment-light than cloud-era software firms: Cursor generates more than twenty times the revenue per employee of cloud-era SaaS comparators.

\paragraph{Contributions.}
This paper makes four contributions.
\begin{itemize}
  \item \textbf{A two-channel model of the AI shock.}
  We embed labor-saving technical change at incumbents and lower fixed costs of entry and operation in a single model of firm dynamics, so that displacement and firm formation are determined jointly rather than studied in isolation.
  \item \textbf{Cloud computing as an empirical analogue.}
  We use the 2006--2020 cloud era in the United States and Europe to discipline the entry-cost channel, accounting for the concurrent decline in business dynamism.
  \item \textbf{Quantitative conditions for offset.}
  We characterize, by calibration and Monte Carlo simulation, the conditions under which employment among firms with 5 to 4,999 employees offsets displacement at large incumbents.
  \item \textbf{Testable empirical signatures.}
  We derive predictions that distinguish offsetting entry from pure displacement and map each one to existing US and European administrative data.
\end{itemize}

The remainder of the paper is organized as follows.
\Cref{sec:theory} presents the framework, \cref{sec:cloud} studies the cloud era, \cref{sec:ai-model} formalizes the AI shock, \cref{sec:simulation} reports the calibration and simulations, \cref{sec:signatures} states the empirical signatures, \cref{sec:discussion} discusses limitations and policy, and \cref{sec:conclusion} concludes.
